\section{总结与延伸}

\subsection{本章小结}

这一讲的结论可以压成一句话：\textbf{数据处理本身就是建模问题，只不过这里的模型要优先考虑吞吐、空间和误差边界，而不是单纯追求精度。}

\begin{importantbox}{三条最终 takeaway}
\begin{enumerate}
    \item KenLM、fastText、DSIR 共享同一类“先打分，再筛选”的过滤范式。
    \item 语言识别、质量过滤、毒性过滤，本质上都是为训练集挑选更合适的样本。
    \item 精确去重、Bloom filter、MinHash 和 LSH 共同构成了大规模语料去重的工具箱。
\end{enumerate}
\end{importantbox}

\begin{knowledgebox}{真正值得记住的不是某个实现}
真正值得记住的是这条线：\textbf{先用粗糙模型估计样本值得不值得留下，再用可扩展的近似算法把这个判断做成百万、亿级数据都能跑完的系统。}
\end{knowledgebox}

\end{document}
